\documentclass[11pt,oneside]{article}

\usepackage[a4paper,margin=26mm,headheight=15pt]{geometry}
\usepackage[T1]{fontenc}
\usepackage[utf8]{inputenc}
\usepackage{libertinus}
\usepackage{microtype}
\usepackage{amsmath,amssymb,amsthm,mathtools,bm,mathrsfs}
\usepackage{booktabs,longtable,array}
\usepackage{graphicx}
\usepackage{xcolor}
\usepackage{enumitem}
\usepackage{fancyhdr}
\usepackage{xurl}
\usepackage{listings}
\usepackage[most]{tcolorbox}
\usepackage{hyperref}
\usepackage[nameinlink,noabbrev]{cleveref}

\definecolor{ProveItBlue}{HTML}{17365D}
\definecolor{ProveItTeal}{HTML}{146C74}
\definecolor{ProveItGold}{HTML}{9B6A00}
\definecolor{ProveItGray}{HTML}{F4F6F8}
\definecolor{CodeGray}{HTML}{F7F7F7}
\definecolor{LinkBlue}{HTML}{0B5CAD}

\hypersetup{
  colorlinks=true,
  linkcolor=ProveItBlue,
  citecolor=ProveItTeal,
  urlcolor=LinkBlue,
  pdftitle={Inverse Asymptotics for a Real-Argument Fibonacci Function},
  pdfauthor={Vladimir Reshetnikov},
  pdfsubject={Log-periodic transseries inversion},
  pdfkeywords={Fibonacci function, inverse asymptotics, transseries, Lagrange inversion, log-periodic expansion}
}

\pagestyle{fancy}
\fancyhf{}
\fancyhead[L]{\small\textsc{ProveIt research note}}
\fancyhead[R]{\small\textsc{Inverse Fibonacci transseries}}
\fancyfoot[C]{\thepage}
\renewcommand{\headrulewidth}{0.4pt}

\setlength{\parindent}{0pt}
\setlength{\parskip}{0.55em}
\setlist{nosep,leftmargin=2em}
\numberwithin{equation}{section}

\newtheorem{theorem}{Theorem}[section]
\newtheorem{proposition}[theorem]{Proposition}
\newtheorem{corollary}[theorem]{Corollary}
\newtheorem{lemma}[theorem]{Lemma}
\newtheorem{definition}[theorem]{Definition}
\newtheorem{remark}[theorem]{Remark}
\newtheorem{example}[theorem]{Example}

\crefname{theorem}{theorem}{theorems}
\crefname{proposition}{proposition}{propositions}
\crefname{corollary}{corollary}{corollaries}
\crefname{lemma}{lemma}{lemmas}
\crefname{definition}{definition}{definitions}
\crefname{remark}{remark}{remarks}
\crefname{example}{example}{examples}

\newtcolorbox{resultbox}[1][]{
  enhanced,breakable,
  colback=ProveItGray,
  colframe=ProveItBlue,
  boxrule=0.7pt,
  arc=1.5mm,
  left=2.5mm,right=2.5mm,top=2mm,bottom=2mm,
  title={#1},
  fonttitle=\bfseries
}
\newtcolorbox{methodbox}[1][]{
  enhanced,breakable,
  colback=white,
  colframe=ProveItTeal,
  boxrule=0.7pt,
  arc=1.5mm,
  left=2.5mm,right=2.5mm,top=2mm,bottom=2mm,
  title={#1},
  fonttitle=\bfseries
}
\newtcolorbox{warningbox}[1][]{
  enhanced,breakable,
  colback=white,
  colframe=ProveItGold,
  boxrule=0.7pt,
  arc=1.5mm,
  left=2.5mm,right=2.5mm,top=2mm,bottom=2mm,
  title={#1},
  fonttitle=\bfseries
}

\lstdefinelanguage{Wolfram}{
  morekeywords={ClearAll,SetDelayed,Module,Table,Sum,Total,Flatten,FullSimplify,
  Assuming,Series,Normal,Expand,Collect,ComplexExpand,Re,Im,Exp,Log,Sqrt,
  Cosh,Cos,Sin,ArcCsch,GoldenRatio,Pi,I,Binomial,N,FindRoot,WorkingPrecision,
  AccuracyGoal,PrecisionGoal,Abs,Max,Plot,NumericQ,Positive,NonNegative,Integers,
  Element},
  sensitive=true,
  morecomment=[s]{(*}{*)},
  morestring=[b]"
}
\lstset{
  language=Wolfram,
  basicstyle=\ttfamily\small,
  keywordstyle=\bfseries\color{ProveItBlue},
  commentstyle=\itshape\color{ProveItTeal},
  stringstyle=\color{ProveItGold},
  backgroundcolor=\color{CodeGray},
  frame=single,
  rulecolor=\color{black!20},
  breaklines=true,
  columns=fullflexible,
  keepspaces=true,
  showstringspaces=false,
  tabsize=2
}

\newcommand{\Fib}{\operatorname{Fib}}
\newcommand{\arsinh}{\operatorname{arsinh}}
\newcommand{\Log}{\operatorname{Log}}
\newcommand{\Res}{\operatorname{Res}}
\newcommand{\val}{\operatorname{val}}
\newcommand{\e}{\mathrm{e}}
\newcommand{\ii}{\mathrm{i}}
\newcommand{\dd}{\mathrm{d}}
\newcommand{\NN}{\mathbb{N}}
\newcommand{\CC}{\mathbb{C}}
\newcommand{\RR}{\mathbb{R}}
\newcommand{\TT}{\mathbb{T}}
\newcommand{\cF}{\mathcal{F}}
\newcommand{\cA}{\mathcal{A}}
\newcommand{\cB}{\mathcal{B}}
\newcommand{\cD}{\mathcal{D}}
\newcommand{\vct}[1]{\boldsymbol{#1}}
\newcommand{\coeff}[2]{[#1]#2}
\newcommand{\bigO}{\mathcal{O}}

\title{\bfseries\LARGE Inverse Asymptotics for a Real-Argument Fibonacci Function\\[0.4em]
\Large Log-periodic transseries, closed all-order coefficients, and a general multiexponential reversion calculus}
\author{Vladimir Reshetnikov\\\small ProveIt research note}
\date{September 3, 2026}

\begin{document}
\maketitle

\begin{abstract}
We study the real interpolation
\[
  \cF(x)=\sqrt{\frac25}\,
  \sqrt{\cosh\!\bigl(2x\,\operatorname{arccsch}2\bigr)-\cos(\pi x)},
  \qquad x\in\RR,
\]
which satisfies \(\cF(n)=\Fib_n\) for every nonnegative integer \(n\).
Because \(\cF\) is even, every sufficiently large value \(y\) has two real preimages,
\(X_\pm(y)=\pm X(y)\).  The dominant inverse is
\(L=\log_\varphi(\sqrt5\,y)\), but the correction is not an ordinary power
series with constant coefficients.  It is a convergent log-periodic transseries.
With
\[
 q=(5y^2)^{-1},\qquad \theta=\pi L,\qquad
 \tau=\frac{\pi}{2\log\varphi},
\]
we prove
\[
 X(y)=L+\frac1{2\log\varphi}\sum_{N\ge1}q^N P_N(\theta;\tau),
\]
where each \(P_N\) is an explicit real trigonometric polynomial containing only
harmonics of the same parity as \(N\).  More strongly, after introducing two
conjugate complex monomials, every coefficient is available in one closed
formula involving generalized binomial coefficients.  The series is the actual
Taylor expansion of a holomorphic implicit solution near the transseries origin,
so total-degree truncation after \(N\) blocks has a uniform
\(\bigO(y^{-2N-2})\) remainder.

The derivation leads to a general inversion calculus for functions of the form
\[
 Y(x)=C\e^{\alpha x}\Psi\!\left(\zeta_1\e^{-\mu_1x},\ldots,
 \zeta_m\e^{-\mu_mx}\right),
 \qquad \Re\mu_j>0,\quad \Psi(\vct0)=1.
\]
For the inverse, the coefficient indexed by a multi-index \(\vct n\ne\vct0\)
is a single diagonal coefficient of the variable power
\(\Psi^{\Gamma_{\vct n}/\alpha}\), where
\(\Gamma_{\vct n}=\vct n\cdot\vct\mu\).  We develop the underlying
Lagrange--B\"urmann identity in filtered transseries algebras, its diagonal
exponential coefficient formula, periodic and quasiperiodic block versions,
conjugation by a nonlinear dominant skeleton, symbolic algorithms, residual
certificates, and numerical validation.
\end{abstract}

\begin{center}
\small
\textbf{Repository context.}
This note is designed as a standalone article for the series-and-transseries
research area of the ProveIt repository:
\url{https://github.com/VladimirReshetnikov/ProveIt/tree/main/Analysis/FabiusFunction/docs/semi-formalized-research-frontiers/drafts/series-and-transseries}.
\end{center}

\clearpage
\begingroup
\small
\setlength{\parskip}{0pt}
\tableofcontents
\endgroup
\clearpage

\section{Problem, notation, and principal result}
\label{sec:intro}

Let
\begin{equation}
  \lambda:=\operatorname{arccsch}2
  =\arsinh\frac12
  =\log\frac{1+\sqrt5}{2}
  =\log\varphi,
  \qquad
  a:=2\lambda,
  \qquad b:=\pi,
  \label{eq:constants}
\end{equation}
where \(\varphi=(1+\sqrt5)/2\) is the golden ratio.  We use
\begin{equation}
  \cF(x):=\sqrt{\frac25}\sqrt{\cosh(ax)-\cos(bx)}.
  \label{eq:F-def}
\end{equation}
The square root is the nonnegative real square root.  The quantities used
throughout the inverse expansion are
\begin{equation}
\begin{aligned}
  R&:=\sqrt5\,y,
  &\qquad
  L&:=\frac{\log R}{\lambda}=\log_\varphi R,\\
  q&:=R^{-2}=\frac1{5y^2}=\e^{-aL},
  &\qquad
  \theta&:=bL,
  \qquad
  \tau:=\frac ba=\frac{\pi}{2\lambda}.
\end{aligned}
  \label{eq:natural-variables}
\end{equation}
Numerically,
\begin{equation}
  \lambda=0.48121182505960344749\ldots,
  \qquad
  \tau=3.26425130263649690673\ldots.
\end{equation}

The dominant estimate \(\cF(x)\sim\varphi^x/\sqrt5\) gives \(x\sim L\).
The oscillatory term \(\cos(\pi x)\), however, survives inversion as a periodic
function of \(\log y\).  It is therefore essential to retain the phase
\(\theta=\pi\log_\varphi(\sqrt5y)\).

For the closed formula introduce the conjugate monomials
\begin{equation}
  \xi:=\e^{-(a+\ii b)L}=q\e^{-\ii\theta}
       =R^{-2-\ii\pi/\lambda},
  \qquad
  \eta:=\e^{-(a-\ii b)L}=q\e^{\ii\theta}
       =R^{-2+\ii\pi/\lambda},
  \label{eq:xi-eta}
\end{equation}
and, for \(r,s\in\NN\),
\begin{equation}
  \chi_{r,s}:=r+s+\ii\tau(r-s).
  \label{eq:chi}
\end{equation}
Here and below \(\NN=\{0,1,2,\ldots\}\), and for \(n\in\NN\),
\(
 \binom zn=z(z-1)\cdots(z-n+1)/n!
\)
with \(\binom z0=1\).

\begin{resultbox}[Principal inverse-transseries formula]
For every real \(y\ge2\), the equation \(\cF(x)=y\) has exactly two real solutions
\(X_\pm(y)=\pm X(y)\), with \(X(y)>0\).  The positive branch has the convergent
expansion
\begin{equation}
 \boxed{
  X(y)=L+\sum_{\substack{r,s\ge0\\r+s\ge1}}
  C_{r,s}\,\xi^r\eta^s,
  \qquad
  C_{r,s}=\frac{(-1)^{r+s+1}}{a\chi_{r,s}}
  \binom{\chi_{r,s}}r\binom{\chi_{r,s}}s .
 }
 \label{eq:main-complex}
\end{equation}
The coefficients obey \(C_{s,r}=\overline{C_{r,s}}\), hence the value on the
physical ray \(\eta=\overline\xi\) is real.  Grouping terms of total degree
\(N=r+s\) gives
\begin{equation}
 \boxed{
  X(y)=L+\frac1a\sum_{N\ge1}q^N P_N(\theta;\tau),
 }
 \label{eq:main-block}
\end{equation}
where
\begin{align}
 P_N(\theta;\tau)
   &:=\sum_{r=0}^{N}\beta_{r,N-r}\,
       \e^{\ii(N-2r)\theta},
 \label{eq:PN-def}\\
 \beta_{r,s}
   &:=aC_{r,s}
    =\frac{(-1)^{r+s+1}}{\chi_{r,s}}
      \binom{\chi_{r,s}}r\binom{\chi_{r,s}}s.
 \label{eq:beta-def}
\end{align}
For each fixed \(N\), \(P_N\) is a real \(2\pi\)-periodic trigonometric
polynomial with harmonics \(N,N-2,\ldots,-N\).  For every fixed truncation
order \(N\),
\begin{equation}
 X(y)=L+\sum_{\substack{r,s\ge0\\1\le r+s\le N}}
 C_{r,s}\xi^r\eta^s+\bigO(q^{N+1})
 =L+\frac1a\sum_{n=1}^{N}q^nP_n(\theta;\tau)
  +\bigO(y^{-2N-2}),
 \label{eq:uniform-remainder}
\end{equation}
uniformly in the phase as \(y\to+\infty\).  The negative branch is obtained
without further calculation:
\begin{equation}
  X_-(y)=-X(y).
\end{equation}
\end{resultbox}

The first four blocks are
\begin{align}
P_1={}&2\cos\theta,
\label{eq:P1}\\
P_2={}&-2-\cos(2\theta)-2\tau\sin(2\theta),
\label{eq:P2}\\
P_3={}&(6-\tau^2)\cos\theta+5\tau\sin\theta
 +\frac{2-9\tau^2}{3}\cos(3\theta)+3\tau\sin(3\theta),
\label{eq:P3}\\
P_4={}&-9+4(3\tau^2-2)\cos(2\theta)
 +\frac{4\tau(2\tau^2-13)}3\sin(2\theta)\notag\\
&\quad+\frac{16\tau^2-1}{2}\cos(4\theta)
 +\frac{\tau(16\tau^2-11)}3\sin(4\theta).
\label{eq:P4}
\end{align}
Thus, in powers of \(y^{-2}\),
\begin{align}
X(y)={}&L+\frac{\cos\theta}{5\lambda y^2}
 -\frac{2+\cos(2\theta)+(\pi/\lambda)\sin(2\theta)}{50\lambda y^4}
\notag\\
&+\frac{1}{250\lambda y^6}
 \left[
 (6-\tau^2)\cos\theta+5\tau\sin\theta
 +\frac{2-9\tau^2}{3}\cos(3\theta)+3\tau\sin(3\theta)
 \right]
 +\bigO(y^{-8}).
\label{eq:main-y-expansion}
\end{align}
The periodic coefficients are not an ornamental rewriting: as shown in
\cref{sec:analytic-status}, no constant-coefficient Poincar\'e series in
integer powers of \(y^{-2}\) can represent this inverse.

\section{The real Fibonacci interpolation and its inverse branches}
\label{sec:branches}

\subsection{Interpolation of the Fibonacci sequence}

The defining expression admits the useful modulus form
\begin{equation}
  5\cF(x)^2
  =\e^{2\lambda x}+\e^{-2\lambda x}-2\cos(\pi x)
  =\left|\e^{\lambda x}-\e^{(-\lambda+\ii\pi)x}\right|^2.
  \label{eq:modulus-form}
\end{equation}
For an integer \(n\ge0\), \(\e^{\ii\pi n}=(-1)^n\), and therefore
\begin{align}
  5\cF(n)^2
  &=\varphi^{2n}+\varphi^{-2n}-2(-1)^n\notag\\
  &=\left(\varphi^n-(-\varphi^{-1})^n\right)^2
   =5\Fib_n^2.
  \label{eq:integer-interpolation}
\end{align}
Both sides are nonnegative, so \(\cF(n)=\Fib_n\).  The function is even and
strictly positive away from the origin.  Near zero,
\begin{equation}
  \cF(x)=\sqrt{\frac{a^2+b^2}{5}}\,|x|+\bigO(|x|^3),
  \label{eq:cusp}
\end{equation}
so the even interpolation has a cusp at \(x=0\), despite the analyticity of
\(\cF(x)^2\).

\subsection{Why only two real branches survive at large height}

Put
\begin{equation}
  h(x):=\cosh(ax)-\cos(bx),
  \qquad \cF(x)^2=\frac25h(x).
\end{equation}
For \(x>0\),
\begin{equation}
  h'(x)=a\sinh(ax)+b\sin(bx).
  \label{eq:h-prime}
\end{equation}
Let
\begin{equation}
  x_*:=\frac1a\arsinh\frac ba.
  \label{eq:x-star}
\end{equation}
For \(x>x_*\), one has \(a\sinh(ax)>b\), hence
\(h'(x)>b+b\sin(bx)\ge0\), with strict positivity.  Define
\begin{equation}
  M:=\max_{0\le x\le x_*}\cF(x).
  \label{eq:M-def}
\end{equation}
Every \(y>M\) has no positive preimage in \([0,x_*]\), while the restriction
of \(\cF\) to \((x_*,\infty)\) is strictly increasing and unbounded.  It has
therefore one positive preimage, and evenness supplies exactly one negative
preimage.

An entirely explicit, if nonoptimal, sufficient height follows from
\(\cos(bx)\ge-1\):
\begin{equation}
  M\le
  \sqrt{\frac25\left(\cosh(ax_*)+1\right)}
  =\sqrt{\frac25\left(\sqrt{1+(b/a)^2}+1\right)}
  =1.3287575092\ldots.
  \label{eq:M-bound}
\end{equation}
For orientation, numerical refinement gives
\(x_*=1.9729927341\ldots\), and the actual maximum on \([0,x_*]\) is
\(1.0138598855\ldots\), attained near \(x=1.1323790680\ldots\).  None of the
asymptotic analysis depends on this numerical refinement.

\begin{remark}[Low-height inverse branches]
The oscillatory term can create several local inverse branches for bounded
values of \(y\).  Those branches terminate at finite critical values and play no
role in the limit \(y\to\infty\).  Throughout the rest of the article,
\(X(y)\) denotes the unique positive outer branch.
\end{remark}

\section{Exact reduction to near-identity reversion}
\label{sec:exact-reduction}

The transseries follows from an exact factorization rather than from repeated
ad hoc substitution.  With
\begin{equation}
  \mu_+:=a+\ii b,
  \qquad
  \mu_-:=a-\ii b,
  \label{eq:mu-pm}
\end{equation}
we have
\begin{align}
  R^2=5y^2
   &=\e^{ax}+\e^{-ax}-2\cos(bx)\notag\\
   &=\e^{ax}\left(1-\e^{-\mu_+x}\right)
                 \left(1-\e^{-\mu_-x}\right).
  \label{eq:factorization}
\end{align}
For real \(x>0\), the two factors are conjugate, and each lies in the open
right half-plane because \(|\e^{-\mu_\pm x}|<1\).  Principal logarithms are
therefore unambiguous.  Taking the positive square root and then the logarithm
gives
\begin{equation}
  \log R=\lambda x+\frac12\log\left(1-\e^{-\mu_+x}\right)
                    +\frac12\log\left(1-\e^{-\mu_-x}\right).
  \label{eq:log-factorization}
\end{equation}
After division by \(\lambda=a/2\), the inverse problem is exactly
\begin{equation}
  L=H(x):=x+A(x),
  \label{eq:H-near-identity}
\end{equation}
where
\begin{align}
  A(x)
  &:=\frac1a\left[
       \log\left(1-\e^{-\mu_+x}\right)
      +\log\left(1-\e^{-\mu_-x}\right)
      \right]
  \label{eq:A-complex}\\
  &=\frac1a\log\left(1-2\e^{-ax}\cos(bx)+\e^{-2ax}\right)
  \label{eq:A-real}\\
  &=-\frac2a\sum_{m\ge1}\frac{\e^{-amx}\cos(mbx)}m.
  \label{eq:A-Fourier}
\end{align}
The last series is absolutely convergent for every \(x>0\).

Set
\begin{equation}
  X(y)=L+U(L).
  \label{eq:X-L-U}
\end{equation}
Then \(U\) is characterized by the exact implicit equation
\begin{equation}
 U+\frac1a\left[
 \log\left(1-\xi\e^{-\mu_+U}\right)
 +\log\left(1-\eta\e^{-\mu_-U}\right)
 \right]=0,
 \label{eq:U-complex-implicit}
\end{equation}
where \(\xi,\eta\) are given by \cref{eq:xi-eta}.  There is also a useful
purely real form.  Dividing the original squared equation by
\(\e^{aL}=R^2\) yields
\begin{equation}
  \e^{aU}+q^2\e^{-aU}-2q\cos\bigl(\theta+bU\bigr)=1.
  \label{eq:U-real-implicit}
\end{equation}
If
\begin{equation}
  W:=aU=\sum_{N\ge1}q^NP_N(\theta;\tau),
  \label{eq:W-def}
\end{equation}
then \cref{eq:U-real-implicit} becomes the compact block-generating equation
\begin{equation}
  \boxed{
  \e^W+q^2\e^{-W}-2q\cos(\theta+\tau W)=1.
  }
  \label{eq:P-generating-equation}
\end{equation}
This identity is an independent way to generate and verify every polynomial
\(P_N\).

\section{Closed all-order coefficients}
\label{sec:closed-coefficients}

\subsection{A one-variable reversion identity}

We begin with the classical near-identity form of the Lagrange--B\"urmann
formula~\cite{Gessel2016,Henrici1974}.  Let \(D=\dd/\dd X\).

\begin{theorem}[Near-identity Lagrange--B\"urmann formula]
\label{thm:near-identity-LB}
Suppose
\begin{equation}
  H(X)=X+A(X)
\end{equation}
has a local inverse \(H^{-1}(X)=X+B(X)\), with \(A\) small in an analytic or
complete filtered formal sense.  Then
\begin{equation}
  \boxed{
  B(X)=\sum_{n\ge1}\frac{(-1)^n}{n!}
  D^{n-1}\!\left(A(X)^n\right).
  }
  \label{eq:LB-near-identity}
\end{equation}
The first terms are
\begin{equation}
  B=-A+AA'-A(A')^2-\frac12A^2A''+\cdots.
  \label{eq:LB-first-terms}
\end{equation}
\end{theorem}

\begin{proof}
Introduce a bookkeeping parameter \(t\) and solve
\begin{equation}
  B_t=-tA(X+B_t),
  \qquad B_0=0.
\end{equation}
The standard Lagrange inversion formula for \(w=t\Phi(X+w)\) gives
\begin{equation}
  w=\sum_{n\ge1}\frac{t^n}{n!}D^{n-1}\!\left(\Phi(X)^n\right).
\end{equation}
Taking \(\Phi=-A\) and then setting \(t=1\) gives
\cref{eq:LB-near-identity}.  Expansion through total degree three in \(A\)
and its derivatives gives \cref{eq:LB-first-terms}.
\end{proof}

For an exponential transseries, differentiation is diagonal on each monomial.
This turns \cref{eq:LB-near-identity} into a particularly economical coefficient
formula.

\begin{proposition}[Diagonal exponential coefficient formula]
\label{prop:diagonal-formula}
Let
\begin{equation}
  A(X)=\sum_{\gamma\in\mathfrak G_+}a_\gamma\e^{-\gamma X},
  \qquad \Re\gamma>0,
  \label{eq:A-exponential-grid}
\end{equation}
where the support is locally finite under addition.  Write
\begin{equation}
  B(X)=\sum_{\gamma\in\mathfrak G_+}b_\gamma\e^{-\gamma X}.
\end{equation}
Then, for every nonzero exponent \(\gamma\),
\begin{align}
  b_\gamma
  &=-\sum_{n\ge1}\frac{\gamma^{n-1}}{n!}
      \coeff{\e^{-\gamma X}}{A(X)^n}
  \label{eq:b-convolution}\\
  &=-\frac1\gamma\coeff{\e^{-\gamma X}}
      {\exp\!\bigl(\gamma A(X)\bigr)}.
  \label{eq:b-diagonal}
\end{align}
Only finitely many values of \(n\) contribute to any fixed coefficient in a
positively filtered grid.
\end{proposition}

\begin{proof}
On the monomial \(\e^{-\gamma X}\),
\begin{equation}
  D^{n-1}\e^{-\gamma X}=(-\gamma)^{n-1}\e^{-\gamma X}.
\end{equation}
Consequently the coefficient contributed by the \(n\)-th summand of
\cref{eq:LB-near-identity} is
\begin{equation}
  \frac{(-1)^n}{n!}(-\gamma)^{n-1}
  \coeff{\e^{-\gamma X}}{A^n}
  =-\frac{\gamma^{n-1}}{n!}
  \coeff{\e^{-\gamma X}}{A^n}.
\end{equation}
Summing over \(n\) proves both forms.
\end{proof}

\begin{remark}[Why the formula is ``diagonal'']
For each target exponent \(\gamma\), one raises the entire perturbation
\(\exp(A)\) to the scalar power \(\gamma\), extracts only the coefficient at
that same exponent, and divides by \(-\gamma\).  No recursive substitution is
needed.  The dependence of the power on the coefficient being extracted is the
feature that packages the inversion combinatorics.
\end{remark}

\subsection{Application to the Fibonacci interpolation}

At \(X=L\), \cref{eq:A-complex} reads
\begin{equation}
  A(L)=\frac1a\bigl(\log(1-\xi)+\log(1-\eta)\bigr).
  \label{eq:A-at-L}
\end{equation}
The derivative \(D=\dd/\dd L\) acts by
\begin{equation}
  D\xi=-\mu_+\xi,
  \qquad
  D\eta=-\mu_-\eta.
  \label{eq:D-xi-eta}
\end{equation}
Thus the monomial \(\xi^r\eta^s\) has exponent
\begin{equation}
  \Gamma_{r,s}:=r\mu_++s\mu_-=a\chi_{r,s}.
  \label{eq:Gamma-rs}
\end{equation}
By \cref{prop:diagonal-formula}, its coefficient in the inverse correction is
\begin{align}
  C_{r,s}
  &=-\frac1{\Gamma_{r,s}}
    \coeff{\xi^r\eta^s}{\exp\!\left(\Gamma_{r,s}A(L)\right)}
  \notag\\
  &=-\frac1{a\chi_{r,s}}
    \coeff{\xi^r\eta^s}{(1-\xi)^{\chi_{r,s}}
                                  (1-\eta)^{\chi_{r,s}}}
  \notag\\
  &=\frac{(-1)^{r+s+1}}{a\chi_{r,s}}
    \binom{\chi_{r,s}}r\binom{\chi_{r,s}}s.
  \label{eq:C-derivation}
\end{align}
This proves the coefficient formula in \cref{eq:main-complex}.  Notice that
\(\Re\chi_{r,s}=r+s>0\), so its denominator never vanishes for
\(r+s\ge1\).

On the physical ray, \(\eta=\overline\xi\).  Since
\begin{equation}
  \chi_{s,r}=\overline{\chi_{r,s}},
  \qquad
  C_{s,r}=\overline{C_{r,s}},
  \label{eq:conjugate-coefficients}
\end{equation}
terms pair to give a real result.  Moreover,
\begin{equation}
  \xi^r\eta^s=q^{r+s}\e^{\ii(s-r)\theta}.
  \label{eq:monomial-block}
\end{equation}
Collecting all pairs with \(r+s=N\) gives \cref{eq:main-block,eq:PN-def}.
The available frequencies are \(s-r=N-2r\), which proves the parity statement
for the harmonics.

An alternative single-line version of the expansion, directly in terms of
complex powers of the large argument, is
\begin{equation}
  X(y)=\frac{\log(\sqrt5y)}\lambda
  +\sum_{\substack{r,s\ge0\\r+s\ge1}}
  C_{r,s}(\sqrt5y)^{-2(r+s)-\ii(\pi/\lambda)(r-s)}.
  \label{eq:complex-power-version}
\end{equation}
Although individual terms are complex, conjugate pairing makes the sum real.
This is the natural generalized-power transseries hidden behind the real
log-periodic blocks.

\section{The log-periodic blocks in detail}
\label{sec:blocks}

\subsection{Generation and parity}

The real logarithmic perturbation at the dominant inverse is
\begin{equation}
  A(L)=-\frac2a\sum_{m\ge1}\frac{q^m\cos(m\theta)}m.
  \label{eq:A-q-theta}
\end{equation}
In the \((q,\theta)\) representation, differentiation with respect to \(L\) is
\begin{equation}
  D=-a q\frac{\partial}{\partial q}
    +b\frac{\partial}{\partial\theta}.
  \label{eq:D-q-theta}
\end{equation}
Substitution of \cref{eq:A-q-theta} into \cref{eq:LB-near-identity}, followed
by collection of powers of \(q\), is a generic way to construct the blocks.
For this particular problem, the binomial formula
\cref{eq:beta-def} is faster and exposes their exact arithmetic structure.
The implicit generating equation \cref{eq:P-generating-equation} gives a third,
independent construction.

For example, the total-degree-one pairs \((r,s)=(1,0),(0,1)\) each have
\(\beta=1\), giving \(P_1=2\cos\theta\).  At total degree two, the central pair
\((1,1)\) contributes \(-2\), while the conjugate edge pair contributes
\(-\cos(2\theta)-2\tau\sin(2\theta)\).  This recovers \cref{eq:P2}.

\begin{proposition}[Block symmetries]
\label{prop:block-symmetries}
For every \(N\ge1\),
\begin{align}
  P_N(\theta+2\pi;\tau)&=P_N(\theta;\tau),
  \label{eq:PN-periodic}\\
  P_N(\theta+\pi;\tau)&=(-1)^N P_N(\theta;\tau),
  \label{eq:PN-antiperiodic}\\
  P_N(-\theta;-\tau)&=P_N(\theta;\tau).
  \label{eq:PN-reflection}
\end{align}
If \(N\) is odd, \(P_N\) has zero phase average.  If \(N=2m\), then
\begin{equation}
  \frac1{2\pi}\int_0^{2\pi}P_{2m}(\theta;\tau)\,\dd\theta
  =-\frac1{2m}\binom{2m}{m}^2.
  \label{eq:PN-average}
\end{equation}
\end{proposition}

\begin{proof}
The first two identities follow from the frequency set
\(N-2r\).  The third follows by complex conjugation.  A constant Fourier term
can occur only when \(r=s=m\), hence only for \(N=2m\).  In that case
\(\chi_{m,m}=2m\), and \cref{eq:beta-def} gives
\begin{equation}
  \beta_{m,m}=-\frac1{2m}\binom{2m}{m}^2.
\end{equation}
\end{proof}

Consequently the phase-averaged inverse correction has the exact series
\begin{equation}
  \left\langle X-L\right\rangle_\theta
  =-\frac1a\sum_{m\ge1}\frac1{2m}\binom{2m}{m}^2q^{2m}.
  \label{eq:mean-correction}
\end{equation}
Using
\begin{equation}
  {}_2F_1\!\left(\frac12,\frac12;1;u\right)
  =\sum_{m\ge0}\binom{2m}{m}^2\frac{u^m}{16^m},
\end{equation}
this can also be written as
\begin{equation}
  \left\langle X-L\right\rangle_\theta
  =-\frac1{2a}\int_0^{16q^2}
  \frac{{}_2F_1(\frac12,\frac12;1;u)-1}{u}\,\dd u.
  \label{eq:mean-hypergeometric}
\end{equation}
The mean is negative, even though the leading pointwise correction
\(q\cos\theta/\lambda\) changes sign.

\subsection{Multiplicative phase recurrence}

Because \(\theta=\pi\log_\varphi(\sqrt5y)\), multiplying \(y\) by
\(\varphi^2\) adds \(2\pi\) to the phase:
\begin{equation}
  \theta(\varphi^2y)=\theta(y)+2\pi,
  \qquad
  q(\varphi^2y)=\varphi^{-4}q(y).
  \label{eq:phase-scaling}
\end{equation}
Thus the coefficient profiles repeat on the multiplicative scale
\(y\mapsto\varphi^2y\), while their amplitudes contract geometrically.
Multiplication by \(\varphi\) adds \(\pi\) to the phase and, by
\cref{eq:PN-antiperiodic}, changes the sign of every odd block.
This discrete scale covariance is the real-variable manifestation of the
complex exponents in \cref{eq:complex-power-version}.

\subsection{Scaled implicit equation as a consistency check}

Writing \(W=\sum_{N\ge1}q^NP_N\), expand
\cref{eq:P-generating-equation} as a formal power series in \(q\).  The
coefficient of \(q\) gives
\begin{equation}
  P_1-2\cos\theta=0.
\end{equation}
At order \(q^2\), one obtains
\begin{equation}
  P_2+\frac12P_1^2+1+2\tau P_1\sin\theta=0,
\end{equation}
which simplifies to \cref{eq:P2}.  Every higher block is determined linearly
after the preceding blocks are known; this makes
\cref{eq:P-generating-equation} particularly effective as a symbolic residual
test.

\section{Analytic status, remainders, and the necessity of transseries}
\label{sec:analytic-status}

\subsection{The expansion is convergent near the transseries origin}

The word \emph{transseries} is sometimes used only for divergent formal
objects; standard background on asymptotic expansions and transseries may be found in~\cite{Olver1997,Edgar2010,vanderHoeven2006,ADH2017}.  In the present problem the generalized expansion is stronger: it is
the Taylor series of a holomorphic function of two independent small
variables.

Define
\begin{equation}
 \mathcal{G}(U,z,w):=
 U+\frac1a\left[
 \log\left(1-z\e^{-\mu_+U}\right)
 +\log\left(1-w\e^{-\mu_-U}\right)
 \right].
 \label{eq:G-holomorphic}
\end{equation}
Near \((U,z,w)=(0,0,0)\), this is holomorphic and satisfies
\begin{equation}
  \mathcal{G}(0,0,0)=0,
  \qquad
  \frac{\partial\mathcal{G}}{\partial U}(0,0,0)=1.
\end{equation}
The holomorphic implicit-function theorem therefore gives a unique convergent
power series
\begin{equation}
  U(z,w)=\sum_{r+s\ge1}C_{r,s}z^rw^s
  \label{eq:U-holomorphic-series}
\end{equation}
in some bidisc \(|z|<\rho\), \(|w|<\rho\).  Formal uniqueness and
\cref{eq:C-derivation} identify its Taylor coefficients with the closed formula
already obtained.

On the physical ray \((z,w)=(\xi,\eta)\),
\begin{equation}
  |z|=|w|=q=\frac1{5y^2}\longrightarrow0.
\end{equation}
Hence \cref{eq:main-complex} converges absolutely for all sufficiently large
\(y\).  Cauchy estimates and coefficient extraction in a smaller bidisc
(compare~\cite{FlajoletSedgewick2009}) give, for every fixed \(N\),
\begin{equation}
 \sum_{r+s\ge N+1}C_{r,s}\xi^r\eta^s=\bigO(q^{N+1}),
 \label{eq:tail-Cauchy}
\end{equation}
with a constant independent of \(\arg\xi=-\theta\).  This proves the uniform
remainder claim \cref{eq:uniform-remainder}.  The same reasoning permits
termwise differentiation with respect to \(L\), and therefore with respect to
\(y\), inside any smaller convergence domain.

The local argument can be supplemented by an explicit sufficient convergence
region.

\begin{proposition}[A concrete bidisc of convergence]
\label{prop:explicit-convergence}
Let
\begin{equation}
  m:=|\mu_\pm|=\sqrt{a^2+b^2},
  \qquad
  k:=\frac{2m}{a}=2\sqrt{1+\tau^2},
\end{equation}
and define
\begin{equation}
  q_{\mathrm{suf}}
  :=\frac{k^k}{(k+1)^{k+1}}
  =0.05024103075902923436\ldots.
  \label{eq:q-sufficient}
\end{equation}
The Taylor series \cref{eq:U-holomorphic-series} converges throughout
\begin{equation}
  \max(|z|,|w|)<q_{\mathrm{suf}}.
  \label{eq:explicit-bidisc}
\end{equation}
Consequently, on the physical ray it converges whenever
\begin{equation}
  y>\frac1{\sqrt{5q_{\mathrm{suf}}}}
   =1.9951967438470964\ldots,
  \label{eq:y-sufficient}
\end{equation}
and in particular for every real \(y\ge2\).
\end{proposition}

\begin{proof}
Rewrite \cref{eq:U-complex-implicit} as \(U=\Phi(U;z,w)\), where
\begin{equation}
 \Phi(U;z,w):=-\frac1a\left[
 \log(1-z\e^{-\mu_+U})+\log(1-w\e^{-\mu_-U})
 \right].
\end{equation}
Choose \(0<\rho<1/(k+1)\) and set
\begin{equation}
 R_\rho:=-\frac2a\log(1-\rho),
 \qquad
 q_\rho:=\rho\e^{-mR_\rho}=\rho(1-\rho)^k.
\end{equation}
If \(|U|\le R_\rho\) and \(|z|,|w|\le q_\rho\), then
\(|z\e^{-\mu_+U}|,|w\e^{-\mu_-U}|\le\rho\).  The elementary bound
\(|\log(1-v)|\le-\log(1-|v|)\) gives
\begin{equation}
  |\Phi(U;z,w)|\le-\frac2a\log(1-\rho)=R_\rho.
\end{equation}
Moreover,
\begin{equation}
  |\partial_U\Phi|
  \le\frac{2m}{a}\frac{\rho}{1-\rho}
  =k\frac{\rho}{1-\rho}<1.
\end{equation}
Thus \(\Phi\) is a uniform contraction of the closed disk
\(|U|\le R_\rho\), and its fixed point depends holomorphically on
\((z,w)\).  The function \(q_\rho=\rho(1-\rho)^k\) increases up to the
endpoint \(\rho=1/(k+1)\), where its supremum is
\(k^k/(k+1)^{k+1}\).  Every smaller bidisc is therefore covered by some
strict contraction, proving the claim.
\end{proof}

\begin{definition}[Block-Poincar\'e expansion]
An expansion
\begin{equation}
  f(y)\sim\sum_{N\ge0}q(y)^N P_N(\theta(y))
\end{equation}
is called a \emph{block-Poincar\'e expansion} if every \(P_N\) is bounded on
its phase domain and truncation after \(N\) has remainder
\(\bigO(q^{N+1})\).  Individual blocks need not form a strict pointwise
asymptotic scale, because a periodic coefficient may vanish at isolated phases.
\end{definition}

The expansion \cref{eq:main-block} is a convergent block-Poincar\'e expansion.
It is also a grid-based exponential transseries in \(L\), and a generalized
complex-power series in \(R=\sqrt5y\).

\subsection{Why an ordinary power series cannot work}

The leading correction is
\begin{equation}
  X(y)-L=\frac q\lambda\cos\theta+\bigO(q^2).
  \label{eq:leading-osc-correction}
\end{equation}
Consider the two sequences
\begin{equation}
 y_n^{(+)}:=\frac{\varphi^{2n}}{\sqrt5},
 \qquad
 y_n^{(-)}:=\frac{\varphi^{2n+1}}{\sqrt5}.
 \label{eq:phase-sequences}
\end{equation}
Along the first, \(\theta=2\pi n\); along the second,
\(\theta=(2n+1)\pi\).  Therefore
\begin{equation}
 \lim_{n\to\infty}\frac{X(y_n^{(+)})-L(y_n^{(+)})}{q(y_n^{(+)})}
 =\frac1\lambda,
 \qquad
 \lim_{n\to\infty}\frac{X(y_n^{(-)})-L(y_n^{(-)})}{q(y_n^{(-)})}
 =-\frac1\lambda.
 \label{eq:no-limit}
\end{equation}
There is no single first constant coefficient.  In particular, no expansion of
the form
\begin{equation}
  X(y)=L+c_1y^{-2}+c_2y^{-4}+\cdots
\end{equation}
can be asymptotic in the usual constant-coefficient sense.  The logarithmic
phase, or equivalently the pair of complex powers, is mathematically forced.

\subsection{Residuals give a posteriori error certificates}

Let \(B_N(L)\) be any approximation to \(U(L)\), and define the exact residual
for the near-identity equation by
\begin{equation}
  \mathcal{R}_N(L):=H\bigl(L+B_N(L)\bigr)-L
  =B_N(L)+A\bigl(L+B_N(L)\bigr).
  \label{eq:residual-def}
\end{equation}
If \(|A'(x)|\le\kappa<1\) on the real segment joining
\(L+B_N(L)\) and \(L+U(L)\), the mean-value theorem gives
\begin{equation}
  |U(L)-B_N(L)|\le\frac{|\mathcal{R}_N(L)|}{1-\kappa}.
  \label{eq:residual-bound}
\end{equation}
Thus a small residual is not merely a numerical diagnostic; once a derivative
bound is supplied it certifies the inverse error.  For the exact transseries
truncation through total degree \(N\), formal substitution gives
\begin{equation}
  \mathcal{R}_N(L)=\bigO(q^{N+1}).
\end{equation}
A Newton correction,
\begin{equation}
  B_N^{\mathrm{new}}
  :=B_N-\frac{\mathcal{R}_N}{1+A'(L+B_N)},
  \label{eq:Newton-correction}
\end{equation}
raises the formal order sharply and is useful when one wants high numerical
accuracy without storing many explicit blocks.

\section{General inversion theory for multiexponential shapes}
\label{sec:general-theory}

The Fibonacci function is one instance of a broad class: a dominant exponential
multiplied by an analytic germ in smaller exponentials.  The same diagonal
coefficient mechanism solves the inverse in closed form.

\subsection{Analytic-germ theorem}

Let \(\vct z=(z_1,\ldots,z_m)\), let
\(\vct n=(n_1,\ldots,n_m)\in\NN^m\), and write
\begin{equation}
  |\vct n|:=n_1+\cdots+n_m,
  \qquad
  \vct z^{\vct n}:=z_1^{n_1}\cdots z_m^{n_m}.
\end{equation}

\begin{theorem}[Inverse of a dominant exponential times an analytic germ]
\label{thm:analytic-germ-inverse}
Let \(\alpha>0\), let \(\mu_1,\ldots,\mu_m\in\CC\) satisfy
\(\Re\mu_j>0\), and let \(\Psi\) be holomorphic and nonzero near
\(\vct0\), with \(\Psi(\vct0)=1\).  Consider a compatible branch of
\begin{equation}
  Y(x)=C\e^{\alpha x}
  \Psi\!\left(\zeta_1\e^{-\mu_1x},\ldots,
              \zeta_m\e^{-\mu_mx}\right),
  \qquad C\ne0.
  \label{eq:general-Y}
\end{equation}
Set
\begin{equation}
  L:=\frac1\alpha\Log\frac YC,
  \qquad
  \xi_j:=\zeta_j\e^{-\mu_jL},
  \qquad
  \Gamma_{\vct n}:=\sum_{j=1}^m n_j\mu_j.
  \label{eq:general-variables}
\end{equation}
Then the inverse branch asymptotic to \(L\) has the convergent multiexponential
expansion
\begin{equation}
  \boxed{
  x=L+\sum_{\vct n\in\NN^m\setminus\{\vct0\}}
       B_{\vct n}\vct\xi^{\vct n},
  \qquad
  B_{\vct n}=-\frac1{\Gamma_{\vct n}}
  \coeff{\vct z^{\vct n}}
  {\Psi(\vct z)^{\Gamma_{\vct n}/\alpha}}.
  }
  \label{eq:general-diagonal-formula}
\end{equation}
The power of \(\Psi\) is defined by the local logarithm satisfying
\(\Log\Psi(\vct0)=0\).  Since \(\Re\Gamma_{\vct n}>0\) for
\(\vct n\ne\vct0\), the denominator does not vanish.
\end{theorem}

\begin{proof}
Taking logarithms in \cref{eq:general-Y} gives
\begin{equation}
  L=x+A(x),
  \qquad
  A(x)=\frac1\alpha\Log\Psi\!\left(
  \zeta_1\e^{-\mu_1x},\ldots,\zeta_m\e^{-\mu_mx}\right).
  \label{eq:general-A}
\end{equation}
At \(x=L\), the monomial \(\vct\xi^{\vct n}\) is an eigenmonomial of
\(D=\dd/\dd L\) with eigenvalue \(-\Gamma_{\vct n}\).  Applying
\cref{prop:diagonal-formula} gives
\begin{align}
 B_{\vct n}
 &=-\frac1{\Gamma_{\vct n}}
   \coeff{\vct z^{\vct n}}
   {\exp\!\left(\Gamma_{\vct n}\alpha^{-1}\Log\Psi(\vct z)\right)}
 \notag\\
 &=-\frac1{\Gamma_{\vct n}}
   \coeff{\vct z^{\vct n}}
   {\Psi(\vct z)^{\Gamma_{\vct n}/\alpha}}.
\end{align}
For convergence, define the analogue of \cref{eq:G-holomorphic}; its derivative
with respect to the inverse correction is \(1\) at the origin.  The
holomorphic implicit-function theorem supplies a convergent multivariate Taylor
series, whose coefficients are uniquely those above.
\end{proof}

\begin{remark}[One-dimensional inversion, multivariate coefficient extraction]
The inverse variable remains one-dimensional.  The several small exponentials
serve only as bookkeeping variables for its transseries.  This is why the
answer requires no multivariate Lagrange determinant: ordinary one-variable
Lagrange inversion, combined with diagonal differentiation of exponential
monomials, already contains the full multi-index combinatorics.
\end{remark}

\subsection{Product-exponential functions}

When the analytic germ factors into binomials, every coefficient in
\cref{eq:general-diagonal-formula} becomes explicit.

\begin{corollary}[Closed formula for product-exponential shapes]
\label{cor:product-exponential}
Under the hypotheses of \cref{thm:analytic-germ-inverse}, suppose
\begin{equation}
  \Psi(\vct z)=\prod_{j=1}^m(1-z_j)^{\beta_j}.
  \label{eq:Psi-product}
\end{equation}
Then
\begin{equation}
  \boxed{
  B_{\vct n}
  =\frac{(-1)^{|\vct n|+1}}{\Gamma_{\vct n}}
   \prod_{j=1}^m
   \binom{\beta_j\Gamma_{\vct n}/\alpha}{n_j}.
  }
  \label{eq:product-coefficients}
\end{equation}
\end{corollary}

\begin{proof}
Raise \cref{eq:Psi-product} to
\(\Gamma_{\vct n}/\alpha\), expand each factor by the generalized binomial
theorem, and use \cref{eq:general-diagonal-formula}.
\end{proof}

For the Fibonacci problem, take
\begin{equation}
  Y=R=\sqrt5\,\cF(x),\quad C=1,\quad \alpha=\lambda,
  \quad (\mu_1,\mu_2)=(a+\ii b,a-\ii b),
  \quad (\beta_1,\beta_2)=\left(\frac12,\frac12\right).
  \label{eq:Fibonacci-general-map}
\end{equation}
Then \(\Gamma_{r,s}/\alpha=2\chi_{r,s}\), so
\(\beta_j\Gamma_{r,s}/\alpha=\chi_{r,s}\), and
\cref{eq:product-coefficients} reduces exactly to
\cref{eq:main-complex}.

\subsection{Dominant exponential polynomials}

A common source of \cref{eq:general-Y} is an exponential polynomial with one
dominant term.  Suppose, for example,
\begin{equation}
  Y(x)^p=c_0\e^{\kappa_0x}
  \left(1+\sum_{j=1}^m c_j\e^{-\nu_jx}\right),
  \qquad \Re\nu_j>0.
  \label{eq:exp-polynomial}
\end{equation}
On a compatible branch,
\begin{equation}
  Y(x)=c_0^{1/p}\e^{(\kappa_0/p)x}
  \left(1+\sum_{j=1}^m c_j\e^{-\nu_jx}\right)^{1/p},
\end{equation}
so \cref{thm:analytic-germ-inverse} applies with
\begin{equation}
  \Psi(\vct z)=\left(1+z_1+\cdots+z_m\right)^{1/p}.
\end{equation}
The Fibonacci square is precisely a three-term exponential polynomial plus a
conjugate oscillatory pair, and the factorization
\cref{eq:factorization} further simplifies its germ into a product.

More generally, any finite exponential sum can be divided by its dominant
monomial.  Taking a logarithm produces an additive semigroup of smaller
exponents, and the diagonal formula applies whether or not the remaining germ
factorizes.

\section{Formal transseries algebra and block calculus}
\label{sec:formal-calculus}

\subsection{A complete exponential-grid algebra}

The preceding derivations can be interpreted analytically, but they also make
sense formally; this filtered viewpoint is consistent with the grid-based and Hahn-style transseries frameworks in~\cite{Edgar2010,vanderHoeven2006,ADH2017}.  This matters when the starting expansion is only asymptotic.
Fix generators \(\mu_1,\ldots,\mu_m\) with positive real parts and form the
additive monoid
\begin{equation}
  \mathfrak G:=\left\{\sum_{j=1}^m n_j\mu_j:\vct n\in\NN^m\right\}.
  \label{eq:exponent-monoid}
\end{equation}
Because the number of generators is finite and \(\Re\mu_j>0\), only finitely
many multi-indices satisfy \(\Re(\vct n\cdot\vct\mu)\le R\) for a fixed
\(R\).  After combining any exact exponent collisions, formal sums
\begin{equation}
  T(X)=\sum_{\gamma\in\mathfrak G}t_\gamma\e^{-\gamma X}
  \label{eq:grid-series}
\end{equation}
form a complete filtered algebra under the real-part filtration.  One may use
\begin{equation}
  \val T:=\min\{\Re\gamma:t_\gamma\ne0\}
  \label{eq:valuation}
\end{equation}
for nonzero \(T\).  Multiplication adds valuations, while differentiation
preserves them:
\begin{equation}
  D\e^{-\gamma X}=-\gamma\e^{-\gamma X}.
\end{equation}
If \(B\) has positive valuation, composition is defined by
\begin{equation}
  \e^{-\gamma(X+B)}=\e^{-\gamma X}\exp(-\gamma B).
  \label{eq:formal-composition}
\end{equation}
Every coefficient below a fixed filtration level receives only finitely many
contributions.

\begin{theorem}[Formal near-identity inversion]
\label{thm:formal-inversion}
Let \(A\) be an exponential-grid series of positive valuation.  There is a
unique positive-valuation series \(B\) satisfying
\begin{equation}
  B(X)=-A(X+B(X)).
  \label{eq:formal-fixed-point}
\end{equation}
It is given by \cref{eq:LB-near-identity}; equivalently, its coefficients are
given by \cref{eq:b-diagonal}.
\end{theorem}

\begin{proof}
Successive approximation in \cref{eq:formal-fixed-point} determines one
filtration layer at a time: replacing \(B\) by another series agreeing through
valuation \(\rho\) changes the right-hand side only at a strictly higher
valuation.  This proves existence and uniqueness in the complete filtered
algebra.  Introducing the formal parameter \(t\) in the proof of
\cref{thm:near-identity-LB} is legitimate coefficientwise, because only
finitely many products contribute at any fixed filtration level.  The same
proof therefore yields the stated formulas.
\end{proof}

This theorem separates two issues that are often conflated:
\begin{enumerate}
  \item formal inversion is controlled by positive valuation and local
        finiteness of the exponent support;
  \item analytic convergence requires additional holomorphy and noncriticality,
        supplied for the Fibonacci problem by the implicit-function theorem.
\end{enumerate}

\subsection{Periodic and quasiperiodic block algebra}

Complex conjugate exponents with a common real part are often clearer after
grouping them into real phase blocks.  Let
\begin{equation}
  q=\e^{-\rho X},
  \qquad
  \vct\theta=\vct\theta_0+\vct\omega X\pmod{2\pi},
  \qquad \rho>0,
  \label{eq:q-torus}
\end{equation}
and suppose
\begin{equation}
  A(X)=\sum_{n\ge1}q^na_n(\vct\theta),
  \label{eq:A-block-general}
\end{equation}
where the \(a_n\) are analytic functions on a torus \(\TT^d\), or formal
Fourier series with appropriate support.  Then
\begin{equation}
  D=-\rho q\partial_q+\vct\omega\cdot\nabla_{\vct\theta},
  \label{eq:D-torus}
\end{equation}
and in particular
\begin{equation}
  D\bigl(q^Nf(\vct\theta)\bigr)
  =q^N\left(\vct\omega\cdot\nabla_{\vct\theta}-\rho N\right)
  f(\vct\theta).
  \label{eq:D-block}
\end{equation}

\begin{theorem}[Universal block reversion formula]
\label{thm:block-reversion}
If
\begin{equation}
  H(X)=X+A(X),
  \qquad
  H^{-1}(X)=X+\sum_{N\ge1}q^Nb_N(\vct\theta),
\end{equation}
then
\begin{equation}
 \boxed{
 b_N(\vct\theta)=
 \sum_{k=1}^{N}\frac{(-1)^k}{k!}
 \left(\vct\omega\cdot\nabla_{\vct\theta}-\rho N\right)^{k-1}
 \sum_{\substack{n_1+\cdots+n_k=N\\n_j\ge1}}
 \prod_{j=1}^k a_{n_j}(\vct\theta).
 }
 \label{eq:block-reversion-formula}
\end{equation}
\end{theorem}

\begin{proof}
The coefficient of \(q^N\) in \(A^k\) is the inner composition sum in
\cref{eq:block-reversion-formula}.  On a block of order \(N\), every
application of \(D\) is the fixed differential operator in
\cref{eq:D-block}.  Extract the coefficient of \(q^N\) from
\cref{eq:LB-near-identity}.
\end{proof}

For the Fibonacci inverse, \(d=1\), \(\rho=a\), \(\omega=b\), and
\begin{equation}
  a_n(\theta)=-\frac{2}{an}\cos(n\theta).
  \label{eq:Fibonacci-a-n}
\end{equation}
The theorem then reproduces \(b_N=P_N/a\).  Unlike the closed binomial formula,
\cref{eq:block-reversion-formula} applies without factorization and is therefore
the natural general-purpose algorithm for log-periodic problems.

If the imaginary parts of the exponent generators have rational rank \(d>1\),
the coefficient functions live on \(\TT^d\) and are quasiperiodic along the
one-dimensional orbit \(X\mapsto\vct\theta_0+\vct\omega X\).  Commensurate
frequencies reduce to an ordinary periodic phase.  Incommensurate decay rates
should usually be left in the full multiexponential grid rather than forced into
a single parameter \(q\).

\subsection{Resonances, reality, and absence of small divisors}

Two distinct multi-indices can occasionally produce the same exponent
\(\Gamma_{\vct n}\).  Their formal monomials must then be combined before one
regards the result as a one-variable exponential series.  The coefficient
formula remains valid termwise in independent bookkeeping variables and the
resonant contributions are subsequently added.

Positive real parts rule out dangerous denominators:
\begin{equation}
  |\Gamma_{\vct n}|\ge\Re\Gamma_{\vct n}>0.
  \label{eq:no-small-divisor}
\end{equation}
Thus the inversion itself creates no small-divisor problem.  For real-valued
functions, exponent generators, germ coefficients, and bookkeeping variables
should be closed under complex conjugation.  The resulting coefficients pair
exactly as in \cref{eq:conjugate-coefficients}.

\subsection{Polynomial and logarithmic coefficient extensions}

If the input contains factors such as \(X^k\e^{-\gamma X}\) or
\((\log X)^k\e^{-\gamma X}\), differentiation is no longer diagonal within a
single monomial.  Polynomial prefactors lead to finite triangular blocks;
logarithms and inverse powers lead to a larger well-ordered differential
coefficient algebra in which differentiation is still triangular.  The
universal formula \cref{eq:LB-near-identity} remains valid unchanged; only the
shortcut \cref{eq:b-diagonal} must be replaced by triangular coefficient
propagation.  This observation is the bridge from pure exponential grids to
logarithmic-exponential transseries.

\section{Conjugation by a nonlinear dominant skeleton}
\label{sec:skeleton}

The leading inverse need not be linear in the logarithm of the target.  A
simple conjugation reduces a much wider class to the same near-identity
problem.

Suppose an inverse equation can be written
\begin{equation}
  Z=S(x)+P(x),
  \label{eq:skeleton-equation}
\end{equation}
where \(S\) is a monotone or locally univalent dominant skeleton with known
inverse \(T=S^{-1}\), and \(P\) is asymptotically small relative to the chosen
scale.  Set
\begin{equation}
  t:=S(x),
  \qquad
  \widetilde A(t):=P(T(t)).
  \label{eq:skeleton-conjugation}
\end{equation}
Then
\begin{equation}
  Z=t+\widetilde A(t),
\end{equation}
so
\begin{equation}
  t=Z+\widetilde B(Z),
  \qquad
  \widetilde B
  =\sum_{n\ge1}\frac{(-1)^n}{n!}
  \frac{\dd^{n-1}}{\dd Z^{n-1}}\left(\widetilde A(Z)^n\right),
  \label{eq:skeleton-B}
\end{equation}
and finally
\begin{equation}
  \boxed{x=T\bigl(Z+\widetilde B(Z)\bigr).}
  \label{eq:skeleton-final}
\end{equation}
If desired, the outer composition is expanded as
\begin{equation}
  x=T(Z)+T'(Z)\widetilde B
  +\frac12T''(Z)\widetilde B^2+\cdots.
  \label{eq:skeleton-outer-expansion}
\end{equation}

This covers dominant skeletons such as
\(S(x)=\alpha x+\beta\log x\), \(S(x)=x\log x\), or a skeleton whose inverse is
expressed by a Lambert \(W\)-function.  The smaller exponential, logarithmic,
and oscillatory corrections are then inverted in the conjugated coordinate
\(t\).  In transseries language, one first inverts the dominant level and only
then reverts the well-based perturbation.

\begin{warningbox}[Scope of the method]
The hypotheses \(\Re\mu_j>0\) and \(\Psi(\vct0)=1\) express genuine decay of
the perturbation after the dominant factor has been removed.  If some
correction has zero real decay, it belongs in the dominant skeleton or in a
phase-dependent leading problem, not in the small transseries.  Singularities
of \(\Psi\), or critical points where the inverse derivative vanishes, determine
the analytic convergence boundary even though the formal expansion remains
well defined.
\end{warningbox}

\section{Algorithms and reproducible coefficient generation}
\label{sec:algorithms}

\subsection{Four complementary algorithms}

The theory suggests four computational routes.

\begin{methodbox}[Algorithm A: direct multi-index coefficients]
For a factored product-exponential input, enumerate the desired multi-indices
\(\vct n\), compute \(\Gamma_{\vct n}\), and evaluate
\cref{eq:product-coefficients}.  This is the fastest route for the Fibonacci
problem and produces exact coefficients in the parameter \(\tau\).
\end{methodbox}

\begin{methodbox}[Algorithm B: analytic-germ diagonal extraction]
For an unfactored germ \(\Psi\), use
\cref{eq:general-diagonal-formula}.  For each \(\vct n\), expand the single
variable power \(\Psi(\vct z)^{\Gamma_{\vct n}/\alpha}\) only far enough to
extract \(\vct z^{\vct n}\).  This avoids constructing irrelevant terms.
\end{methodbox}

\begin{methodbox}[Algorithm C: universal block formula]
When decay orders are commensurate, represent the perturbation as
\(A=\sum q^na_n(\vct\theta)\) and apply
\cref{eq:block-reversion-formula}.  Fourier or polynomial representations of
the coefficient functions make the phase differential operator sparse.
\end{methodbox}

\begin{methodbox}[Algorithm D: triangular implicit recurrence]
Use the fixed-point equation
\begin{equation}
  B=-A(X+B)
   =-\sum_{m\ge0}\frac{B^m}{m!}D^mA(X).
  \label{eq:triangular-recurrence}
\end{equation}
At order \(N\), the right-hand side involves only coefficients of \(B\) from
lower orders, because both \(A\) and \(B\) have positive valuation.  This gives
a memory-efficient recursive implementation.  Substitution into the original
implicit equation provides an immediate residual check.
\end{methodbox}

For the Fibonacci blocks, direct expansion of
\cref{eq:P-generating-equation} is especially convenient.  If
\(W_{<N}=\sum_{n=1}^{N-1}q^nP_n\), then the coefficient of \(q^N\) in
\begin{equation}
  \e^{W_{<N}+q^NP_N}+q^2\e^{-W_{<N}-q^NP_N}
  -2q\cos\bigl(\theta+\tau(W_{<N}+q^NP_N)\bigr)-1
\end{equation}
is linear in \(P_N\); setting it to zero determines the next block.

\subsection{Ordering and truncation}

For a finite grid with unequal decay rates, total multi-index degree is not
always the best notion of order.  A robust implementation uses the weight
\begin{equation}
  w(\vct n):=\Re\Gamma_{\vct n}
  \label{eq:weight-order}
\end{equation}
and retains all terms below a prescribed weight.  In the Fibonacci problem,
all generators have real part \(a\), so
\(w(r,s)=a(r+s)\) and total degree is exactly the decay order.

When terms are grouped into phase blocks, the remainder should be stated in
block norm, as in \cref{eq:uniform-remainder}, rather than by dividing by a
possibly vanishing trigonometric polynomial.  This convention keeps error
statements uniform and numerically meaningful.

\subsection{Symbolic verification protocol}

A reliable calculation should perform all three checks below:
\begin{enumerate}
  \item verify conjugation, \(C_{s,r}=\overline{C_{r,s}}\), and realness of each
        grouped block;
  \item substitute the truncated \(W\) into
        \cref{eq:P-generating-equation} and confirm that coefficients through
        the requested order vanish;
  \item compare the truncated inverse against a high-precision numerical root
        of \(\cF(x)=y\), preferably at both generic values and exact Fibonacci
        values.
\end{enumerate}
The Wolfram Language implementation in \cref{app:WL} follows this protocol.

\section{Numerical validation}
\label{sec:numerics}

For \(N\ge0\), define the truncation
\begin{equation}
  X_N(y):=L+\frac1a\sum_{n=1}^{N}q^nP_n(\theta;\tau),
  \qquad X_0(y):=L.
  \label{eq:XN-def}
\end{equation}
The positive reference root was computed at high precision from
\(\cF(x)=y\), starting from \(L\).  The table reports absolute errors.
The coefficients through \(P_6\) are given in
\cref{sec:intro,app:higher-blocks}.

\begin{table}[htbp]
\centering
\caption{Absolute inverse errors \(E_N(y)=|X_N(y)-X(y)|\).}
\label{tab:numerical-errors}
\resizebox{\textwidth}{!}{%
\begin{tabular}{@{}r l rrrrr@{}}
\toprule
\(y\) & exact or numerical \(X(y)\) & \(E_0\) & \(E_1\) & \(E_2\) & \(E_4\) & \(E_6\)\\
\midrule
2   & \(3\)                    & \(1.12696\times10^{-1}\) & \(1.52360\times10^{-2}\) & \(2.96187\times10^{-3}\) & \(3.76916\times10^{-4}\) & \(1.28682\times10^{-4}\)\\
3   & \(4\)                    & \(4.47122\times10^{-2}\) & \(1.01264\times10^{-3}\) & \(4.22181\times10^{-4}\) & \(8.66820\times10^{-6}\) & \(2.17452\times10^{-7}\)\\
5   & \(5\)                    & \(1.68278\times10^{-2}\) & \(2.26345\times10^{-4}\) & \(1.85965\times10^{-5}\) & \(2.98868\times10^{-8}\) & \(1.82331\times10^{-11}\)\\
10  & \(6.4577930902882116495\) & \(5.45185\times10^{-4}\) & \(1.13522\times10^{-5}\) & \(1.54252\times10^{-7}\) & \(3.01709\times10^{-11}\) & \(6.00018\times10^{-15}\)\\
100 & \(11.242189770132803894\) & \(3.01016\times10^{-5}\) & \(3.56196\times10^{-9}\) & \(2.97159\times10^{-13}\) & \(4.30138\times10^{-21}\) & \(1.48686\times10^{-28}\)\\
\bottomrule
\end{tabular}%
}
\end{table}

The first three target values are Fibonacci numbers:
\(2=\Fib_3\), \(3=\Fib_4\), and \(5=\Fib_5\), so their exact positive
preimages are the integers shown.  The table illustrates two features.
First, the dominant logarithm \(L\) is already useful.  Second, retaining the
phase-dependent blocks produces the predicted rapid improvement as
\(q=(5y^2)^{-1}\) decreases.  At \(y=100\), six blocks give about twenty-eight
decimal digits in the inverse.

\subsection{Exact consistency at Fibonacci targets}

Binet's formula gives, for \(n\ge1\),
\begin{equation}
  \sqrt5\,\Fib_n
  =\varphi^n\left(1-(-1)^n\varphi^{-2n}\right).
  \label{eq:Binet-factor}
\end{equation}
Hence, when \(y=\Fib_n\),
\begin{equation}
  L=n+\frac1\lambda\log\left(1-(-1)^n\varphi^{-2n}\right),
  \label{eq:L-at-Fib}
\end{equation}
and the exact inverse correction is
\begin{equation}
  X(y)-L=n-L
  =-\frac1\lambda\log\left(1-(-1)^n\varphi^{-2n}\right).
  \label{eq:exact-correction-at-Fib}
\end{equation}
The convergent transseries must sum to this quantity at every sufficiently
large Fibonacci target.  This supplies a discrete infinite family of exact
checks in addition to generic numerical root comparisons.

\subsection{Derivative of the inverse}

Termwise differentiation of \cref{eq:main-block} gives asymptotics for the
sensitivity of the inverse.  From the first block,
\begin{equation}
  \frac{\dd X}{\dd L}
  =1-\frac q\lambda\left(a\cos\theta+b\sin\theta\right)+\bigO(q^2).
\end{equation}
Since \(\dd L/\dd y=(\lambda y)^{-1}\),
\begin{equation}
  \boxed{
  \frac{\dd X}{\dd y}
  =\frac1{\lambda y}\left[
   1-q\left(2\cos\theta+\frac\pi\lambda\sin\theta\right)
   +\bigO(q^2)
  \right].
  }
  \label{eq:inverse-derivative}
\end{equation}
This agrees with the implicit identity \(X'(y)=1/\cF'(X(y))\) on the outer
branch.

\section{Structural conclusions and extensions}
\label{sec:consequences}

The calculation reveals several general principles.

\paragraph{Complex exponents are the simplest coordinates.}
The real oscillation \(\cos(\pi x)\) becomes a conjugate pair of decaying
exponentials after the dominant term is factored.  In those coordinates,
reversion is diagonal and the coefficient formula is one line.  Trigonometric
blocks should be formed only after inversion.

\paragraph{Log-periodicity is inherited from the ratio of exponents.}
The decay rate is \(a=2\log\varphi\) and the angular frequency is \(b=\pi\).
Their ratio \(\tau=b/a\) controls every sine coefficient.  In a general
conjugate pair \(\rho\pm\ii\omega\), the corresponding ratio
\(\omega/\rho\) plays the same role.

\paragraph{A finite factorization can encode infinitely many inverse terms.}
The exact factorization \cref{eq:factorization} contains only two small factors,
yet generalized binomial extraction generates the complete two-dimensional
coefficient array \(C_{r,s}\).  No term-by-term matching of the original
inverse equation is required.

\paragraph{Formal and analytic conclusions should be separated.}
Positive valuation proves formal invertibility.  Holomorphic implicit-function
theory proves convergence.  A problem may possess the first property without
the second, especially when the supplied direct expansion is merely asymptotic.
The same coefficient identities still determine the formal inverse.

\paragraph{Nonlinear leading behavior is handled by conjugation.}
When the dominant inverse is not a logarithm, one first inverts the skeleton
\(S\), then applies near-identity reversion in the skeleton coordinate, and
finally composes with \(S^{-1}\).  This organizes mixed logarithmic,
exponential, and oscillatory scales without changing the core inversion
identity.

\paragraph{Further questions.}
The exact bivariate Taylor series invites a study of its singular variety and
therefore of the sharp convergence domain in \((\xi,\eta)\).  Other natural
directions are asymptotics of \(C_{r,s}\) along rays in the index lattice,
complex continuation of the inverse branches, and automated support management
for general Hahn-style transseries with polynomial-logarithmic coefficients.

\section{Conclusion}
\label{sec:conclusion}

For the real-argument Fibonacci interpolation \cref{eq:F-def}, the large real
inverse is governed by the dominant logarithm
\(L=\log_\varphi(\sqrt5y)\), but its correction is intrinsically
log-periodic.  The exact answer is the convergent conjugate-pair transseries
\cref{eq:main-complex}, with every coefficient given by
\begin{equation*}
  C_{r,s}=\frac{(-1)^{r+s+1}}{2\log\varphi}\,
  \frac{\binom{r+s+\ii\tau(r-s)}r
        \binom{r+s+\ii\tau(r-s)}s}
       {r+s+\ii\tau(r-s)}.
\end{equation*}
Its real form is a series in \((5y^2)^{-1}\) whose coefficients are finite
Fourier polynomials in \(\pi\log_\varphi(\sqrt5y)\).  Total-degree truncation
has a uniform algebraic remainder, and the negative inverse branch is simply
the negative of the positive one.

The same derivation yields the general diagonal formula
\cref{eq:general-diagonal-formula} for a dominant exponential multiplied by an
analytic germ in smaller exponentials.  Product factorizations reduce it to
generalized binomial coefficients; commensurate complex exponents produce
periodic blocks; incommensurate frequencies produce torus-valued quasiperiodic
blocks; and a nonlinear dominant skeleton is handled by conjugation.  These
formulas are consequences and reorganizations of classical
Lagrange--B\"urmann inversion, presented here in a form adapted to exponential
and logarithmic transseries.  No claim of priority is made for the general
classical inversion identities; the contribution of this note is the explicit
Fibonacci inverse, its closed coefficient array, and the unified specialization
to multiexponential shapes.

\appendix

\section{Higher log-periodic blocks}
\label{app:higher-blocks}

For reference, the next two blocks generated by \cref{eq:beta-def} are
\begin{align}
P_5={}&
 \frac{\tau^4-95\tau^2+240}{6}\cos\theta
 +\frac{4\tau(31-2\tau^2)}{3}\sin\theta
\notag\\
&+\frac{27\tau^4-213\tau^2+40}{4}\cos(3\theta)
 -\frac{7\tau(9\tau^2-11)}{2}\sin(3\theta)
\notag\\
&+\frac{625\tau^4-875\tau^2+24}{60}\cos(5\theta)
 -\frac{25\tau(5\tau^2-1)}{6}\sin(5\theta),
\label{eq:P5}\\[0.4em]
P_6={}&-\frac{200}{3}
 -\frac{92\tau^4-937\tau^2+450}{6}\cos(2\theta)
 -\frac{\tau(8\tau^4-418\tau^2+1035)}{6}\sin(2\theta)
\notag\\
&-\frac{4(64\tau^4-116\tau^2+9)}{3}\cos(4\theta)
 -\frac{4\tau(64\tau^4-620\tau^2+261)}{15}\sin(4\theta)
\notag\\
&-\frac{324\tau^4-135\tau^2+2}{6}\cos(6\theta)
 -\frac{\tau(648\tau^4-1530\tau^2+137)}{30}\sin(6\theta).
\label{eq:P6}
\end{align}
The constant in \(P_6\) agrees with
\cref{eq:PN-average}:
\begin{equation}
  -\frac16\binom63^2=-\frac{200}{3}.
\end{equation}

For exact symbolic computation, an equivalent rising-factorial form of the
complex coefficient is often convenient:
\begin{equation}
  C_{r,s}
  =-\frac{(-\chi_{r,s})^{\overline r}
           (-\chi_{r,s})^{\overline s}}
          {a\,r!\,s!\,\chi_{r,s}},
  \label{eq:C-rising-factorial}
\end{equation}
where \(z^{\overline n}=z(z+1)\cdots(z+n-1)\).  This avoids gamma-function
representations whose removable singularities can obscure the fact that the
generalized binomial coefficient is a polynomial for integer lower index.

\section{Wolfram Language implementation}
\label{app:WL}

The following code evaluates the closed transseries directly.  It keeps the
complex conjugate form internally and takes the real part only after summation.

\begin{lstlisting}
ClearAll["Global`*"];

lam = Log[GoldenRatio];
aa = 2 lam;
tau0 = Pi/aa;

chi[r_Integer?NonNegative, s_Integer?NonNegative] /; r + s > 0 :=
  r + s + I tau0 (r - s);

cc[r_Integer?NonNegative, s_Integer?NonNegative] /; r + s > 0 :=
  (-1)^(r + s + 1) Binomial[chi[r, s], r] Binomial[chi[r, s], s]/
    (aa chi[r, s]);

ell[y_?NumericQ] /; y > 0 := Log[Sqrt[5] y]/lam;
qq[y_?NumericQ] /; y > 0 := 1/(5 y^2);
theta[y_?NumericQ] /; y > 0 := Pi ell[y];

xApprox[y_?NumericQ, nmax_Integer?NonNegative, wp_: 50] /; y > 0 :=
 Module[{yy = N[y, wp]},
  N[Re[
    ell[yy] +
     Sum[
      cc[r, n - r] qq[yy]^n Exp[I (n - 2 r) theta[yy]],
      {n, 1, nmax}, {r, 0, n}]
    ], wp]
  ];

f[x_?NumericQ] :=
  Sqrt[2/5] Sqrt[Cosh[2 lam x] - Cos[Pi x]];

xExact[y_?NumericQ, wp_: 80] /; y > 0 :=
 Module[{xx, yy = N[y, wp]},
  xx /. FindRoot[
    f[xx] == yy,
    {xx, N[ell[yy], wp]},
    WorkingPrecision -> wp,
    AccuracyGoal -> Floor[wp/2],
    PrecisionGoal -> Floor[wp/2]]
  ];

errorTable[ys_List, orders_List, wp_: 80] :=
 Table[
  With[{xe = xExact[y, wp]},
   Prepend[Abs[xApprox[y, #, wp] - xe] & /@ orders, xe]
   ],
  {y, ys}
  ];

errorTable[{2, 3, 5, 10, 100}, {0, 1, 2, 4, 6}]
\end{lstlisting}

For exact symbolic Fourier blocks, it is useful to keep \(\tau\) as an
independent symbol and to implement the generalized binomial coefficient as a
finite product.

\begin{lstlisting}
ClearAll[gbin, chiT, betaT, p, z, t];

gbin[u_, 0] := 1;
gbin[u_, n_Integer?Positive] :=
  Product[u - k, {k, 0, n - 1}]/n!;

chiT[r_Integer?NonNegative, s_Integer?NonNegative, t_] /;
   r + s > 0 := r + s + I t (r - s);

betaT[r_Integer?NonNegative, s_Integer?NonNegative, t_] /;
   r + s > 0 :=
 Module[{ch = chiT[r, s, t]},
  (-1)^(r + s + 1) gbin[ch, r] gbin[ch, s]/ch
  ];

p[n_Integer?Positive, z_, t_] :=
 FullSimplify[
  TrigReduce@ComplexExpand[
    Sum[
     betaT[r, n - r, t] Exp[I (n - 2 r) z],
     {r, 0, n}]
    ],
  Assumptions -> Element[{z, t}, Reals]
  ];

Table[p[n, z, t], {n, 1, 6}]
\end{lstlisting}

Finally, the real generating equation provides a residual test independent of
the complex coefficient construction.

\begin{lstlisting}
ClearAll[wTrunc, blockResidual, qsym];

wTrunc[nmax_Integer?NonNegative, z_, qsym_, t_] :=
  Sum[qsym^n p[n, z, t], {n, 1, nmax}];

blockResidual[nmax_Integer?NonNegative, z_, qsym_, t_] :=
 Module[{w = wTrunc[nmax, z, qsym, t]},
  FullSimplify[
   Normal@Series[
     Exp[w] + qsym^2 Exp[-w] -
      2 qsym Cos[z + t w] - 1,
     {qsym, 0, nmax}],
   Assumptions -> Element[{z, t}, Reals]
   ]
  ];

(* This returns 0: all coefficients through qsym^6 vanish. *)
blockResidual[6, z, qsym, t]
\end{lstlisting}

\section{A compact general-purpose derivation template}
\label{app:template}

For reuse on another inverse problem of similar shape, the derivation can be
organized as follows.

\begin{enumerate}
  \item \textbf{Factor the dominant term.}  Rewrite the direct function as
  \begin{equation}
    Y(x)=C\e^{\alpha x}\Psi(\vct\zeta\e^{-\vct\mu x}),
    \qquad \Psi(\vct0)=1.
  \end{equation}

  \item \textbf{Choose the dominant inverse.}  Set
  \begin{equation}
    L=\alpha^{-1}\Log(Y/C).
  \end{equation}

  \item \textbf{Introduce independent small monomials.}  Put
  \begin{equation}
    \xi_j=\zeta_j\e^{-\mu_jL}.
  \end{equation}
  Keep conjugate monomials independent during coefficient extraction.

  \item \textbf{Apply the diagonal formula.}  For every retained multi-index,
  calculate
  \begin{equation}
    B_{\vct n}=-\Gamma_{\vct n}^{-1}
    \coeff{\vct z^{\vct n}}
      {\Psi(\vct z)^{\Gamma_{\vct n}/\alpha}}.
  \end{equation}

  \item \textbf{Group equal decay orders and conjugates.}  Convert complex
  monomials into periodic or quasiperiodic real coefficient blocks only after
  their coefficients are known.

  \item \textbf{State the remainder in the correct norm.}  If the analytic
  implicit-function theorem applies, use a uniform block remainder.  If the
  direct expansion is formal, state the result in the corresponding filtered
  transseries algebra.

  \item \textbf{Verify.}  Substitute the truncation into the exact or truncated
  defining equation, inspect the first nonzero residual order, and compare with
  high-precision numerical inversion where available.
\end{enumerate}

\begin{thebibliography}{99}

\bibitem{Gessel2016}
I.~M. Gessel,
\newblock Lagrange inversion,
\newblock \emph{Journal of Combinatorial Theory, Series A} \textbf{144}
(2016), 212--249.

\bibitem{Henrici1974}
P. Henrici,
\newblock \emph{Applied and Computational Complex Analysis, Volume 1},
\newblock Wiley, New York, 1974.

\bibitem{FlajoletSedgewick2009}
P. Flajolet and R. Sedgewick,
\newblock \emph{Analytic Combinatorics},
\newblock Cambridge University Press, Cambridge, 2009.

\bibitem{Olver1997}
F.~W.~J. Olver,
\newblock \emph{Asymptotics and Special Functions},
\newblock A K Peters, Wellesley, 1997.

\bibitem{Edgar2010}
G.~A. Edgar,
\newblock Transseries for beginners,
\newblock \emph{Real Analysis Exchange} \textbf{35} (2010), 253--310.

\bibitem{vanderHoeven2006}
J. van der Hoeven,
\newblock \emph{Transseries and Real Differential Algebra},
\newblock Lecture Notes in Mathematics 1888, Springer, Berlin, 2006.

\bibitem{ADH2017}
M. Aschenbrenner, L. van den Dries, and J. van der Hoeven,
\newblock \emph{Asymptotic Differential Algebra and Model Theory of
Transseries},
\newblock Annals of Mathematics Studies 195, Princeton University Press,
Princeton, 2017.

\end{thebibliography}

\end{document}
